% ═══════════════════════════════════════════════════════════════
% ML-04b: El horario (MUSTERLÖSUNG)
% Spanisch Klasse 7 - Sequenz 4: Mi colegio
% ═══════════════════════════════════════════════════════════════
% Identisches Layout wie AB-04b, Lösungen in ROT
% ═══════════════════════════════════════════════════════════════

\documentclass[11pt, a4paper]{article}

% ═══════════════════════════════════════════════════════════════
% SW-MODUS TOGGLE
% ═══════════════════════════════════════════════════════════════
\newif\ifswmodus
\swmodusfalse

% ═══════════════════════════════════════════════════════════════
% PAKETE
% ═══════════════════════════════════════════════════════════════

\usepackage[utf8]{inputenc}
\usepackage[T1]{fontenc}
\usepackage[ngerman]{babel}
\usepackage{helvet}
\renewcommand{\familydefault}{\sfdefault}

\usepackage[margin=2cm]{geometry}
\usepackage{enumitem}
\usepackage{tabularx}
\usepackage{booktabs}
\usepackage{multirow}
\usepackage{array}

\usepackage{xcolor}
\usepackage{framed}
\usepackage{tcolorbox}
\tcbuselibrary{skins}

\usepackage{fancyhdr}
\usepackage{lastpage}
\usepackage{needspace}

% ═══════════════════════════════════════════════════════════════
% FARBEN
% ═══════════════════════════════════════════════════════════════

\definecolor{spanisch-color}{HTML}{c41e3a}
\definecolor{grau-color}{HTML}{888888}
\definecolor{hellgrau-color}{HTML}{f5f5f5}
\definecolor{orange-color}{HTML}{b35c00}
\definecolor{loesung-color}{HTML}{c62828}

\definecolor{sw-dark}{HTML}{000000}
\definecolor{sw-gray}{HTML}{666666}
\definecolor{sw-white}{HTML}{ffffff}

\ifswmodus
    \colorlet{spanisch}{sw-dark}
    \colorlet{grau}{sw-gray}
    \colorlet{hellgrau}{sw-white}
    \colorlet{orange}{sw-dark}
    \colorlet{loesung}{sw-dark}
\else
    \colorlet{spanisch}{spanisch-color}
    \colorlet{grau}{grau-color}
    \colorlet{hellgrau}{hellgrau-color}
    \colorlet{orange}{orange-color}
    \colorlet{loesung}{loesung-color}
\fi

\newcommand{\fachfarbe}{spanisch}

% ═══════════════════════════════════════════════════════════════
% VARIABLEN
% ═══════════════════════════════════════════════════════════════

\newcommand{\thema}{El horario -- Stundenplan und Uhrzeiten}
\newcommand{\sequenz}{04}
\newcommand{\block}{b}
\newcommand{\fach}{Spanisch}
\newcommand{\klasse}{7}
\newcommand{\maxpunkte}{20}
\newcommand{\datum}{\today}

% ═══════════════════════════════════════════════════════════════
% KOPF- UND FUSSZEILE
% ═══════════════════════════════════════════════════════════════

\pagestyle{fancy}
\fancyhf{}
\renewcommand{\headrulewidth}{0.4pt}
\renewcommand{\footrulewidth}{0.4pt}

\lhead{\fach{} -- Klasse \klasse}
\chead{\textcolor{loesung}{\textbf{ML-\sequenz\block{} (MUSTERLÖSUNG)}}}
\rhead{\datum}

\lfoot{}
\cfoot{}
\rfoot{Seite \thepage\ von \pageref{LastPage}}

% ═══════════════════════════════════════════════════════════════
% BEFEHLE
% ═══════════════════════════════════════════════════════════════

\newcommand{\punktebox}[1]{\hfill\fbox{\small \_/#1}}
\newcommand{\luecke}[1]{\rule{#1}{0.4pt}}
\newcommand{\lsg}[1]{\textcolor{loesung}{\textbf{#1}}}

\newtcolorbox{merkkasten}[1][]{
    colback=hellgrau,
    colframe=\fachfarbe,
    colbacktitle=white,
    coltitle=\fachfarbe,
    boxrule=1pt,
    arc=0pt,
    left=8pt, right=8pt, top=8pt, bottom=8pt,
    fonttitle=\bfseries,
    attach boxed title to top left={yshift=-2mm, xshift=5mm},
    boxed title style={boxrule=0pt, colback=white},
    title=#1
}

\newtcolorbox{vokabelkasten}{
    colback=white,
    colframe=\fachfarbe,
    boxrule=0.5pt,
    arc=0pt,
    left=5pt, right=5pt, top=5pt, bottom=5pt
}

\newtcolorbox{hinweis}{
    colback=white,
    colframe=orange,
    colbacktitle=white,
    coltitle=orange,
    boxrule=1pt,
    arc=0pt,
    left=5pt, right=5pt, top=5pt, bottom=5pt,
    fonttitle=\bfseries,
    attach boxed title to top left={yshift=-2mm, xshift=5mm},
    boxed title style={boxrule=0pt, colback=white},
    title=Hinweis
}

\newenvironment{aufgabe}[1]{%
    \needspace{5\baselineskip}%
    \nopagebreak[4]%
    \vspace{10pt}
    \noindent\textbf{\textcolor{\fachfarbe}{Aufgabe #1}}
    \nopagebreak[4]%
    \vspace{5pt}
    \nopagebreak[4]%
    \begin{list}{}{\leftmargin=0pt}
    \item[]
}{%
    \end{list}
}

\newcommand{\es}[1]{\textcolor{\fachfarbe}{\textbf{#1}}}

% ═══════════════════════════════════════════════════════════════
% DOKUMENT
% ═══════════════════════════════════════════════════════════════

\begin{document}

% ═══════════════════════════════════════════════════════════════
% KOPFBEREICH
% ═══════════════════════════════════════════════════════════════

\begin{center}
    {\Large\bfseries\textcolor{\fachfarbe}{Arbeitsblatt: \thema}}\\[0.3em]
    {\large\textcolor{loesung}{\textbf{--- MUSTERLÖSUNG ---}}}
\end{center}

\vspace{5pt}

\noindent
\begin{tabularx}{\textwidth}{|X|X|X|}
\hline
\textbf{Name:} \lsg{Musterschüler/in} &
\textbf{Klasse:} \klasse &
\textbf{Datum:} \datum \\
\hline
\end{tabularx}

\vspace{15pt}

% ═══════════════════════════════════════════════════════════════
% MERKKASTEN: Die Wochentage
% ═══════════════════════════════════════════════════════════════

\begin{merkkasten}[Merkkasten: Los días de la semana (Die Wochentage)]
\begin{center}
\begin{tabular}{|c|c|c|c|c|c|c|}
\hline
\es{L} & \es{M} & \es{X} & \es{J} & \es{V} & \es{S} & \es{D} \\
\hline
lunes & martes & miércoles & jueves & viernes & sábado & domingo \\
\hline
\lsg{Montag} & \lsg{Dienstag} & \lsg{Mittwoch} & \lsg{Donnerstag} & \lsg{Freitag} & \lsg{Samstag} & \lsg{Sonntag} \\
\hline
\end{tabular}
\end{center}

\vspace{0.3em}
\textbf{Frage:} \es{¿Qué día es hoy?} = Welcher Tag ist heute?
\end{merkkasten}

\vspace{10pt}

% ═══════════════════════════════════════════════════════════════
% MERKKASTEN: Uhrzeiten
% ═══════════════════════════════════════════════════════════════

\begin{merkkasten}[Merkkasten: Las horas (Die Uhrzeiten)]
\begin{tabular}{ll}
\es{¿Qué hora es?} & Wie spät ist es? \\[0.3em]
\es{Es la una.} & Es ist ein Uhr. \\
\es{Son las dos / tres / ...} & Es ist zwei / drei / ... Uhr. \\[0.3em]
\es{Son las ocho y media.} & Es ist halb \lsg{neun}. \\
\es{a las nueve} & um neun Uhr \\
\es{de ocho a nueve} & von acht bis neun \\
\end{tabular}

\vspace{0.5em}
\textbf{Frage:} \es{¿A qué hora tienes...?} = Um wie viel Uhr hast du...?
\end{merkkasten}

\vspace{10pt}

% ═══════════════════════════════════════════════════════════════
% AUFGABE 1: Wochentage ordnen
% ═══════════════════════════════════════════════════════════════

\begin{aufgabe}{1 -- Ordne die Wochentage}
Bringe die Wochentage in die richtige Reihenfolge (1-7). \punktebox{4}
\end{aufgabe}

\begin{vokabelkasten}
domingo | martes | viernes | lunes | sábado | jueves | miércoles
\end{vokabelkasten}

\vspace{0.5em}

\noindent
\begin{tabular}{ccccccc}
\fbox{\phantom{00}1\phantom{00}} & \fbox{\phantom{00}2\phantom{00}} & \fbox{\phantom{00}3\phantom{00}} & \fbox{\phantom{00}4\phantom{00}} & \fbox{\phantom{00}5\phantom{00}} & \fbox{\phantom{00}6\phantom{00}} & \fbox{\phantom{00}7\phantom{00}} \\[0.5em]
\lsg{lunes} & \lsg{martes} & \lsg{miércoles} & \lsg{jueves} & \lsg{viernes} & \lsg{sábado} & \lsg{domingo} \\
\end{tabular}

\vspace{15pt}

% ═══════════════════════════════════════════════════════════════
% AUFGABE 2: Uhrzeiten schreiben
% ═══════════════════════════════════════════════════════════════

\begin{aufgabe}{2 -- Schreibe die Uhrzeiten auf Spanisch}
Schreibe die Uhrzeiten in Worten. \punktebox{4}
\end{aufgabe}

\noindent
\begin{tabular}{p{3cm}p{10cm}}
a) 8:00 & \lsg{Son las ocho.} \\[1em]
b) 10:30 & \lsg{Son las diez y media.} \\[1em]
c) 1:00 & \lsg{Es la una.} \\[1em]
d) 3:30 & \lsg{Son las tres y media.} \\
\end{tabular}

\vspace{15pt}

% ═══════════════════════════════════════════════════════════════
% AUFGABE 3: Stundenplan ausfüllen
% ═══════════════════════════════════════════════════════════════

\begin{aufgabe}{3 -- Completa tu horario (Fülle deinen Stundenplan aus)}
Trage drei Schulfächer mit Tag und Uhrzeit ein. \punktebox{6}
\end{aufgabe}

\begin{hinweis}
Schulfächer: \es{matemáticas}, \es{español}, \es{inglés}, \es{educación física}, \es{música}, \es{arte}
\end{hinweis}

\vspace{0.5em}

\noindent
\begin{tabularx}{\textwidth}{|X|X|X|}
\hline
\textbf{Asignatura (Fach)} & \textbf{Día (Tag)} & \textbf{Hora (Uhrzeit)} \\
\hline
\lsg{matemáticas} & \lsg{lunes} & \lsg{a las ocho} \\[0.5em]
\hline
\lsg{español} & \lsg{martes} & \lsg{a las diez} \\[0.5em]
\hline
\lsg{educación física} & \lsg{viernes} & \lsg{a las once} \\[0.5em]
\hline
\end{tabularx}

\textit{(Beispiellösung -- individuelle Antworten möglich)}

\vspace{15pt}

% ═══════════════════════════════════════════════════════════════
% AUFGABE 4: Dialog schreiben
% ═══════════════════════════════════════════════════════════════

\begin{aufgabe}{4 -- Escribe un diálogo (Schreibe einen Dialog)}
A fragt B nach dem Stundenplan. Verwende die Fragen aus dem Merkkasten. \punktebox{6}
\end{aufgabe}

\begin{hinweis}
Verwende: \es{¿Cuándo tienes...?} (Wann hast du...?) und \es{¿A qué hora...?} (Um wie viel Uhr...?)
\end{hinweis}

\vspace{0.5em}

\noindent
\textbf{A:} \lsg{¿Cuándo tienes matemáticas?}

\vspace{1em}

\noindent
\textbf{B:} \lsg{Los lunes a las ocho.}

\vspace{1em}

\noindent
\textbf{A:} \lsg{¿Y español?}

\vspace{1em}

\noindent
\textbf{B:} \lsg{Los martes a las diez.}

\vspace{1em}

\noindent
\textbf{A:} \lsg{¿A qué hora tienes educación física?}

\vspace{1em}

\noindent
\textbf{B:} \lsg{Los viernes a las once.}

\textit{(Beispiellösung -- Bewertung: Korrekte Fragestrukturen 3P, passende Antworten 3P)}

% ═══════════════════════════════════════════════════════════════
% PUNKTEÜBERSICHT
% ═══════════════════════════════════════════════════════════════

\vspace{20pt}
\hrule
\vspace{10pt}

\begin{flushright}
\textbf{Punkte:} \lsg{\maxpunkte} / \maxpunkte
\end{flushright}

\end{document}
